\section{Complete nonlinear maps: symmetry, balance and periods}
\label{sec:nonlinear}

The arguments in this section use exactness on the complete finite state
space.  They do not linearize a decoder on a correctness source.  A complete
bijection has balanced column outputs; an arbitrary function agreeing with
transpose on a subspace need not have them.  This distinction supplies an
additional obstruction beyond the source-dimension bounds of Part I.

\subsection{One nonuniform layer and periodic first maps}

\begin{theorem}[An equivariant sandwich]\label{thm:uniform-sandwich}
\formal{UniformSandwich.lean}{uniform_sandwich_injective_classes,uniform_sandwich_finite}
Let the entry alphabet have at least two elements.  Suppose $U,V$ commute
with every column rotation, $U$ is injective, $K$ is class-uniform, and
$UKV=\tau$ on the complete state space.  Then every receiver class of $K$
is a singleton.  If the alphabet is finite, the injectivity of $U$ follows
from the factorization and need not be assumed.
\end{theorem}
\begin{outline}
Put $S=\operatorname{Fix}(T)$ and $W=KV(S)$.  The set $\tau S$ consists
of matrices with all rows equal, so it is invariant under column rotation.
Since $UW=\tau S$, equivariance and injectivity of $U$ imply that $W$ is
column-rotation invariant.  On each input in $V(S)\subseteq S$, the columns
of $K$ agree at any two receivers $c,e$ in the same class.  Applying this
fact to every rotation of every element of $W$ shows that $T^{e-c}$ fixes
all of $W$.  Equivariance of $U$ would then make $T^{e-c}$ fix all of
$\tau S$.  An arbitrary repeated row word disproves this when $c\ne e$.
For a finite alphabet, $U$ is surjective because the composite is, and
therefore injective.
\end{outline}

The following quantitative class count concerns an actual periodic first
map, and permits that map to be nonlinear.

\begin{theorem}[A periodic first factor]\label{thm:periodic-first}
\formal{SubgroupClasses.lean}{covariance_no_class_collision,orbit_classes_injective,exact_periodic_first_class_bound,exact_periodic_first_class_size_bound}
Let $A$ be a finite alphabet with at least two elements and let
$\tau=KF$ on $A^{d\times d}$, where $F$ is any total map commuting
with $T^r$ and $K$ is class-uniform with $M$ classes.  Put
$g=\gcd(d,r)$.  Every class meets each orbit of addition by $r$ in at
most one receiver.  Consequently each class has at most $g$ receivers and
\begin{equation}\label{eq:periodic-class-bound}
 M\ge d/g.
\end{equation}
The conclusion includes $r=0$, when it says only $M\ge1$.
\end{theorem}
\begin{outline}
Both factors are bijective.  From $FT^r=T^rF$ and $KF=\tau$ we obtain
$KT^r=R^rK$, hence the same covariance for every $s\in\langle r\rangle$.
If $e=c+s$ lies in the same class as $c$, the shared rule and this covariance
give, for every input $X$,
\[
 \operatorname{col}_eK(X)
 =\operatorname{col}_cK(T^sX)
 =\sigma^s\operatorname{col}_cK(X).
\]
Surjectivity of $K$ rules out this fixed relation between two distinct
output columns.  The $g$ cyclic orbits give the asserted class-size and
class-count bounds.  The formal proof derives the numerical bound through
the order of $r$ in $\Z/d\Z$.
\end{outline}

The bound is sharp for every divisor $g\mid d$.  Here is the explicit
partial-digit-swap construction, a reuse of the composition examples in
\cite{SeventeenthPaper2026}.  Write $d=gh$, $i=gw+z$ and $c=gu+v$, with
$0\le z,v<g$ and $0\le w,u<h$.  Define
\[
 (FX)_{gw+z,gu+v}=X_{gw+v,gu+z},\qquad
 (KY)_{gw+z,gu+v}=Y_{gu+z,gw+v}.
\]
Both are coordinate permutations, $F$ commutes with $T^g$, and direct
substitution gives $KF=\tau$.  The second map has the $h$ interval classes
$p(gu+v)=u$, with
\[
 f_u(Z)_{gw+z}=Z_{gu+z,\,g(w-u)\bmod d}.
\]
Indeed $Z=T^{gu+v}Y$ turns the column index in this formula into $gw+v$.
If $\gcd(r,d)=g$, then $\langle r\rangle=\langle g\rangle$, so this
also realizes equality for the prescribed subgroup.  The construction
is algebraic; a coordinate permutation is not a free tape operation.

\subsection{Agreement distance from equivariance}

The one-layer obstruction has a simple probabilistic form on an arbitrary
correctness set.  This form uses no linear-source assumption.

\begin{proposition}[Full-input agreement bound]\label{prop:agreement}
\formal{CompositionAgreement.lean}{full_input_agreement_probability,repeated_class_agreement_probability}
Let $A$ be finite, $|A|\ge2$, and put $q=|A|^d$, the cardinality of one
whole column.  If $K$ is class-uniform, with class sizes $n_J$, then
\begin{equation}\label{eq:agreement-bound}
 \Pr_X[KX=\tau X]
 \le \frac{M}{d}+
       \frac{\sum_J n_J(n_J-1)}{2dq},
\end{equation}
where $X$ is uniform on the complete state space.  If some class has two
distinct receivers, the separate bound
\begin{equation}\label{eq:agreement-pair}
 \Pr_X[KX=\tau X]\le\frac{1+q^{-1}}2
\end{equation}
also holds.
\end{proposition}
\begin{outline}
For a fixed $X$, consider all $d$ rotated inputs $T^sX$ and count the
phases on which $K$ is correct.  At receiver $c=-s$, its view is $X$.
Correctness there requires
\[
 f_{p(c)}(X)=z_c(X),\qquad
 z_c(X):=\sigma^{-c}\operatorname{row}_cX.
\]
The $z_c$ are independent uniform words from a set of cardinality $q$.
In one class, the number of correct phases is at most the largest
multiplicity of these words, hence at most one plus the number of equal
unordered pairs in that class.  Taking expectations and using rotation
invariance of the uniform input law proves \cref{eq:agreement-bound}.

For the second bound, take distinct same-class receivers $c,e$, set
$s=e-c$, and let $G=\{X:KX=\tau X\}$.  If both $X$ and $T^sX$ lie in
$G$, sharedness gives
$\operatorname{row}_eX=\sigma^s\operatorname{row}_cX$.
This event has probability $1/q$, since the two rows are independent.
Thus $\Pr(G\cap T^{-s}G)\le1/q$.  Inclusion--exclusion and equality of
the two marginal probabilities give $2\Pr(G)-1\le1/q$.
\end{outline}

For example, let $F$ be any first map, let $H$ be equivariant, and let
$\gamma=\Pr[KF\ne\tau]$ and $\epsilon=\Pr[F\ne H]$.  The composite
$KH$ has the same class partition as $K$, and is correct on every input
where $KF$ is correct and $F=H$.  Therefore
$\gamma+\epsilon$ is at least one minus either applicable bound above.
This measures equality of entire maps on an input.  It does not measure
the number of output coordinates changed by the first map.

\subsection{Balanced intertwiners}

We now allow a second nonuniform layer, but retain a concrete property
of its first nonuniform operator.  Let $\Sigma$ be a finite alphabet of
size $a\ge2$ with a distinguished symbol $0$.  For width $b\ge1$, regard
a state as a tuple of $db$ words in $\Sigma^d$, one for each row and lane.
Let $L$ be a zero-fixing permutation of $\Sigma^d$, applied identically
to all these words.  Define
\[
 Q=LTL^{-1},\qquad Q_c=T^cQT^{-c}.
\]
In the next statement the group acts on scalar words; the remaining
maps may mix all rows and lanes.

\begin{theorem}[Transitive pairs obstruct balanced outputs]
\label{thm:balanced-pair}
\formal{CompositionWordAction.lean}{transitive_pair_impossible}
Suppose $H$ is a total column-equivariant map and $K$ is a total
class-uniform map on $\Sigma^{d\times d\times b}$, where $d\ge2$.
If $KLH=\tau$ everywhere and $p(c)=p(e)$, the group
$\langle Q_c,Q_e\rangle$ cannot act transitively on
$\Sigma^d\setminus\{0\}$.
\end{theorem}
\begin{outline}
Exactness on the finite complete domain makes $H$ and $K$ bijective.
Writing $q=a^{db}$ for the whole-column alphabet, each column of $K$
has exactly $q^{d-1}$ preimages for each value.  Since input rotation is
bijective, every shared decoder $f_J$ has these same balanced fibres.

Equivariance of $H$ gives $KQ=RK$.  Evaluating at a receiver in $J$ and
changing its view yields
\begin{equation}\label{eq:decoder-covariance}
 f_J(Q_cY)=\sigma f_J(Y)\qquad(c\in J).
\end{equation}
The corresponding inverse identities also hold.  A word in $Q_c,Q_e$
and their inverses therefore rotates $f_J$ by its signed exponent sum.
If the word represents the identity in the actual scalar-word group
$G=\langle Q_c,Q_e\rangle$, it fixes every matrix.  Surjectivity of
$f_J$ and the exact order $d$ of $\sigma$ imply that its exponent sum
is zero modulo $d$.  Consequently there is a surjective homomorphism
\[
 \theta:G\longrightarrow\Z/d\Z,
 \qquad\theta(Q_c)=\theta(Q_e)=1.
\]
Let $N=\ker\theta$.  Equation~\eqref{eq:decoder-covariance} makes $f_J$
constant on all $N$-orbits of matrices.

Suppose $G$ were transitive on the nonzero words.  Normality of $N$
makes all its orbits there have one size $\lambda$, and $G/N$ acts
transitively on the set of these orbits.  Its cardinality $t$ divides
$d$, so
\[
 \lambda=(a^d-1)/t>1,\qquad \gcd(\lambda,a)=1.
\]
The strict inequality follows from $t\le d<a^d-1$ for $a,d\ge2$.
If a matrix $Y$ is nonzero, choose one nonzero row/lane word $y$.
The stabilizers satisfy $N_Y\le N_y$, hence
\[
 |N\cdot Y|=[N:N_Y]=[N:N_y][N_y:N_Y]
            =\lambda[N_y:N_Y].
\]
This uses no transitivity on tuples of words.  Choose a column value
different from $f_J(0)$.  Its fibre excludes the zero matrix and is a
union of $N$-orbits, so its size is divisible by $\lambda$.  But its
balanced size $q^{d-1}$ is a power of $a$, a contradiction.
\end{outline}

Surjectivity alone would give the cyclic quotient but would not give
the last divisibility contradiction.  Likewise, a code decoder defined
on a proper correctness source need not have balanced fibres.

\begin{proposition}[Anchored binary differences]\label{prop:anchored}
\formal{CompositionAnchoredEndpoint.lean}{anchored_adjacent_classes_distinct}
For $d\ge3$, define the binary word permutation
\[
 (Lx)_0=x_0,\qquad (Lx)_i=x_i+x_{i+1}\quad(1\le i<d),
\]
with cyclic indices.  In any complete factorization $KLH=\tau$, with
$H$ equivariant, the partition of $K$ satisfies $p(c)\ne p(c+1)$ for
every receiver $c$.  This holds at every positive width, even when
$H$ and $K$ are nonlinear and mix lanes.
\end{proposition}
\begin{outline}
The inverse is $x_0=y_0$ and
$x_i=y_0+\sum_{j=i}^{d-1}y_j$ for $1\le i<d$.
For $S=\sum_i y_i$, this gives
$Qy=(S,y_2,\ldots,y_{d-1},S+y_0)$.
Conjugating $\langle Q,TQT^{-1}\rangle$ by $L^{-1}$ gives
$\langle T,A\rangle$, where direct substitution in this inverse yields
\[
 Ax=(x_2,x_1+x_2,x_0+x_1,x_3+\eta,\ldots,x_{d-1}+\eta),
 \qquad\eta=x_0+x_1+x_2.
\]
On its first triple, $A$ has the cycle
\[
 100\to001\to110\to010\to011\to101\to111\to100.
\]
Moreover $(Ax)_j+(Ax)_2=x_j+x_2$ for every $j\ge3$.
Thus a power taking a nonzero triple to another with the same third
bit fixes every outside coordinate, regardless of its value.  In
particular the following are local moves in the group:
\[
 100\to010,\quad110\to010,\quad101\to001,
 \quad011\to001,\quad111\to001.
\]
They use powers $3,1,3,4,2$, respectively.  Conjugation by rotations
places them on any cyclic triple.

For a word with at least two ones, choose consecutive ones in cyclic
order.  If their intervening zero gap has length at least two, use
$100\to010$ to shorten it.  A gap of length zero or one then permits
one of the other moves to lower the weight.  Iteration reaches a unit
vector, which a rotation sends to coordinate zero.  Reversing such
paths proves transitivity on all nonzero words, including when $d=3$.
Conjugacy gives the same transitivity for every adjacent pair
$\langle Q_c,Q_{c+1}\rangle$.  Apply \cref{thm:balanced-pair}.
\end{outline}

\subsection{A joint-output period constraint}

Balance examines a shared output column.  A literal period of the final
mask supplies a separate constraint on the complete output map.

\begin{theorem}[Transport of a mask period]\label{thm:period-transport}
\formal{PeriodTransport.lean}{literal_period_complete,fixed_conjugate_iff,fixed_image_preserved}
Let the entry alphabet be finite.  Suppose $U,B,H$ are total
column-equivariant maps, $L$ is any total map, $K$ is class-uniform,
and
\[
 \tau=UKBLH
\]
on the complete state space.  Exactness makes $L$ bijective.  If
$p(c+s)=p(c)$ for every $c$, then its actual conjugated shift
$Q=LTL^{-1}$ commutes with $T^s$.
\end{theorem}
\begin{outline}
All five factors are bijective.  Put $F=UKB$.  Sharedness and the
literal mask period give $KT^s=T^sK$, hence $FT^s=T^sF$.
Equivariance of $H$ gives $FQ=RF$.  As $R$ commutes with $T^s$,
\[
 FQT^s=RFT^s=RT^sF=T^sRF=T^sFQ=FT^sQ.
\]
Cancel the injective map $F$.  The formal proof derives the required
equivalences and injectivity from the total factorization.
\end{outline}

In particular, $T^s$ must preserve
$\operatorname{Fix}(Q)=L\operatorname{Fix}(T)$.
For the anchored binary map of \cref{prop:anchored}, this is the space
of matrices supported on column zero: $L(1,\ldots,1)=e_0$ on each
row/lane word.  No nonzero column rotation preserves this space.

\begin{corollary}[The anchored two-class architecture]\label{cor:anchored-all}
\formal{CompositionAnchoredEndpoint.lean}{anchored_at_most_two_classes_impossible}
For every $d\ge3$ and every positive binary width, there is no complete
factorization $\tau=KBLH$ in which $L$ is the anchored difference,
$B,H$ are arbitrary equivariant total maps, and $K$ uses at most two
receiver classes.
\end{corollary}
\begin{outline}
Absorb $B$ into $K$; this preserves its shared partition.
By \cref{prop:anchored}, that partition is a proper coloring of the
cycle.  An odd cycle has no proper two-coloring.  On an even cycle,
every proper two-coloring alternates and has period two.  This is a
nonzero period for even $d\ge4$, contradicting
\cref{thm:period-transport} and the preceding fixed-space calculation.
\end{outline}

The period theorem permits an arbitrary equivariant outer $U$.  The
combined anchored corollary does not: its balance step concerns $K$
itself.  Neither result excludes two arbitrary nonlinear nonuniform
layers, gives a positive agreement-rate deficit for the anchored
architecture, or supplies a native running-time bound.
